\subsection{Product of a finite number of measures}

Let \(X_{i}\) ( \(1 \leqslant i \leqslant n\) ) be \(n\) locally compact spaces, \(X = \prod_{i=1}^{n} X_{i}\) their product. The set of linear combinations of complex functions of the form

\[
\left(x _ {1}, x _ {2}, \dots , x _ {n}\right) \mapsto f _ {1} \left(x _ {1}\right) f _ {2} \left(x _ {2}\right) \dots f _ {n} \left(x _ {n}\right),
\]

where \(f_{i} \in \mathcal{K}(\mathrm{X}_{i};\mathbf{C})\) ( \(1 \leqslant i \leqslant n\) ), may be identified canonically with the tensor product \(\bigotimes_{i=1}^{n} \mathcal{K}(\mathrm{X}_{i};\mathbf{C})\) , and it follows from Lemma 1 of No.~1, by induction on \(n\) , that this tensor product is dense in \(\mathcal{K}(\mathrm{X};\mathbf{C})\) .

Now let \(\mu_{i}\) be a measure on \(X_{i}\) ( \(1 \leqslant i \leqslant n\) ); there exists on X one and only one measure \(\nu\) such that, for \(f_{i} \in \mathcal{K}(X_{i}; \mathbf{C})\) ( \(1 \leqslant i \leqslant n\) ),

\[
\left\langle f _ {1} \otimes f _ {2} \otimes \dots \otimes f _ {n}, \nu \right\rangle = \prod_ {i = 1} ^ {n} \left\langle f _ {i}, \mu_ {i} \right\rangle .\tag{11}
\]

For, if this measure exists, it is unique by the foregoing. On the other hand, let \(\nu = \mu_{1} \otimes \mu_{2} \otimes \cdots \otimes \mu_{n}\) be the measure on X defined by the recursion relation

\[
\mu_ {1} \otimes \mu_ {2} \otimes \dots \otimes \mu_ {n} = (\mu_ {1} \otimes \mu_ {2} \otimes \dots \otimes \mu_ {n - 1}) \otimes \mu_ {n}.
\]

It follows from No.~1, formula (1) and this definition (by induction on n) that \(\nu\) verifies (11); it is said to be the product measure of the measures \(\mu_{i}\) ( \(1 \leqslant i \leqslant n\) ) and it is denoted again by \(\bigotimes_{i=1}^{n} \mu_{i}\) . The relation (11) may also be written

\[
\left\langle f _ {1} \otimes f _ {2} \otimes \dots \otimes f _ {n}, \mu_ {1} \otimes \mu_ {2} \otimes \dots \otimes \mu_ {n} \right\rangle = \prod_ {i = 1} ^ {n} \left\langle f _ {i}, \mu_ {i} \right\rangle .\tag{12}
\]

PROPOSITION 7 (`associativity of the product of measures'). --- Let \((\mathbf{I}_k)_{1 \leqslant k \leqslant r}\) be a partition of the interval \([1, n]\) of \(\mathbf{N}\) ; then

\[
\bigotimes_ {k = 1} ^ {r} \left(\bigotimes_ {i \in I _ {k}} \mu_ {i}\right) = \bigotimes_ {i = 1} ^ {n} \mu_ {i}\tag{13}
\]

For, these two measures coincide, by (12), for every function in \(\bigotimes_ {i = 1} ^ {n} \mathcal{K}(X_{i}; \mathbf{C})\) .

The integral of a function \(f \in \mathcal{K}(\mathbf{X};\mathbf{C})\) with respect to the product measure is denoted

\[
\int f d \mu_ {1} d \mu_ {2} \dots d \mu_ {n},
\]

or

\[
\iint \ldots \int f d \mu_ {1} d \mu_ {2} \ldots d \mu_ {n}
\]

or

\[
\int f (\mu_ {1} \otimes \dots \otimes \mu_ {n})
\]

or also

\[
\iint \dots \int f (x _ {1}, x _ {2}, \dots , x _ {n}) d \mu_ {1} (x _ {1}) d \mu_ {2} (x _ {2}) \dots d \mu_ {n} (x _ {n})
\]

or

\[
\iint \dots \int f (x _ {1}, x _ {2}, \dots , x _ {n}) \mu_ {1} (x _ {1}) \mu_ {2} (x _ {2}) \dots \mu_ {n} (x _ {n})
\]

with n signs \(\int\) ; it is said to be a multiple integral of order n, or an n-tuple integral. By virtue of the associativity of the product of measures and the theorem on inverting the order of integration (No.~1, Th. 2), we have, for every permutation \(\sigma\) of \([1,n]\) ,

\[
\iint \dots \int f d \mu_ {1} d \mu_ {2} \dots d \mu_ {n} = \int d \mu_ {\sigma (1)} \int d \mu_ {\sigma (2)} \dots \int f d \mu_ {\sigma (n)}.\tag{14}
\]

The integral notation and formula (14) may be extended in an obvious way to functions \(\mathbf{f} \in \mathcal{K}(\mathrm{X};\mathrm{E})\) with values in a Hausdorff locally convex space E, such that \(\mathbf{f}(\mathrm{X})\) is contained in a complete convex subset of E. We leave to the reader the task of generalizing to the product of any finite number of measures the results of Nos. 2 and 3 concerning the product of two measures.

In particular, one calls Lebesgue measure on \(\mathbf{R}^n\) the product of \(n\) measures identical to the Lebesgue measure on \(\mathbf{R}\) ; the integral of a function \(\mathbf{f} \in \mathcal{K}(\mathbf{R}^n; \mathbf{E})\) , satisfying the preceding condition, is denoted

\[
\iint \dots \int \mathbf {f} (x _ {1}, x _ {2}, \dots , x _ {n}) d x _ {1} d x _ {2} \dots d x _ {n}
\]

and is equal to

\[
\int_ {- \infty} ^ {+ \infty} d x _ {1} \int_ {- \infty} ^ {+ \infty} d x _ {2} \dots \int_ {- \infty} ^ {+ \infty} \mathbf {f} (x _ {1}, x _ {2}, \dots , x _ {n}) d x _ {n}.
\]

Lebesgue measure on \(R^{n}\) is invariant under every translation.
% semantic-fragments: legacy-input-seam
\relax{}
